% ============================================================
% UESTC Physics Experiments I — POSTLAB Template
% REQUIRES: XeLaTeX (NOT pdfLaTeX)
% ============================================================
% Fill in YOUR info and data below, then compile twice:
%   xelatex -interaction=nonstopmode file.tex
%   xelatex -interaction=nonstopmode file.tex
% ============================================================
\documentclass[12pt,a4paper]{article}
\usepackage{pdfpages}
\usepackage[a4paper,top=1in,bottom=1in,left=1.25in,right=1.25in,headsep=10pt]{geometry}
\usepackage{iftex}
\ifPDFTeX
  \usepackage{newtxtext,newtxmath}
  \newcommand{\tablefont}{\sffamily}
  \newcommand{\titlefont}{\rmfamily}
\else
  \usepackage{fontspec}
  \setmainfont{Times New Roman}
  \setsansfont{Times New Roman}
  \setmonofont{Courier New}
  % === Choose your table font strategy ===
  % Strategy A (clean, uniform — recommended):
  \newfontfamily\tablefont{Times New Roman}
  \newfontfamily\titlefont{Times New Roman}
  \newfontfamily\watermarkfont{Times New Roman}
  % Strategy B (Calibri tables — uncomment and comment out Strategy A):
  % \IfFontExistsTF{Calibri}{\newfontfamily\tablefont{Calibri}}{\newfontfamily\tablefont{Arial}}
  % \IfFontExistsTF{Calibri Light}{\newfontfamily\watermarkfont{Calibri Light}}{\newfontfamily\watermarkfont{Arial}}
  % \newfontfamily\titlefont{Times New Roman}
  \usepackage{newtxmath}
\fi
\usepackage{amsmath}
\usepackage{array,booktabs,tabularx,longtable,float,graphicx,caption}
\usepackage{enumitem}
\usepackage{fancyhdr}
\usepackage{xcolor}
\usepackage{eso-pic}
\usepackage{tikz}
\usetikzlibrary{calc}
\usepackage{setspace}
\usepackage{indentfirst}

% ---- Page layout (adjust stretch for content density) ----
\setlength{\parindent}{2em}
\setlength{\parskip}{1pt}                           % 0.5pt for compact
\setstretch{1.04}                                    % 1.02 for compact/long reports

% ---- Math display spacing ----
\setlength{\abovedisplayskip}{5pt plus 1pt minus 2pt}
\setlength{\belowdisplayskip}{5pt plus 1pt minus 2pt}
\setlength{\abovedisplayshortskip}{4pt plus 1pt minus 2pt}
\setlength{\belowdisplayshortskip}{4pt plus 1pt minus 2pt}

% ---- Float spacing (uncomment for compact layout) ----
% \setlength{\textfloatsep}{6pt plus 1pt minus 2pt}
% \setlength{\floatsep}{5pt plus 1pt minus 2pt}
% \setlength{\intextsep}{5pt plus 1pt minus 2pt}

\captionsetup{font=small,labelfont=bf,skip=4pt}      % skip=2pt for compact

% ---- Page style ----
\pagestyle{fancy}
\fancyhf{}
\fancyfoot[C]{\fontsize{11pt}{11pt}\selectfont \thepage}
\renewcommand{\headrulewidth}{0.5pt}
\renewcommand{\footrulewidth}{0pt}
\setlength{\headheight}{14pt}
\setcounter{page}{1}

% ============================================================
% === FILL IN YOUR INFORMATION HERE ===
% ============================================================
\newcommand{\CoverPDF}{02-Template for lab report-2026.pdf}
\newcommand{\DataPDF}{scanned-data.pdf}
% Add figure filenames as needed:
% \newcommand{\FigOne}{fig_example.pdf}
% \newcommand{\FigTwo}{fig_example2.pdf}

% ============================================================
% === SECTION & TABLE MACROS (do not edit) ===
% ============================================================

% Score box in left margin
\newcommand{\scorebox}{%
  \begin{tabular}{|>{\centering\arraybackslash}m{16mm}|}
  \hline
  \textbf{Score}\\[9mm]
  \hline
  \end{tabular}%
}

% Section heading with score box — standard spacing
% For compact: change vspace before to 0.35em, after to 0.4em
\newcommand{\labsection}[2]{%
  \par\vspace{0.5em}
  \noindent
  \begin{minipage}[t]{16mm}
    \vspace{0pt}
    \scorebox
  \end{minipage}\hspace{6mm}%
  \begin{minipage}[t]{\dimexpr\textwidth-22mm\relax}
    \vspace{2mm}
    {\bfseries\fontsize{14pt}{16pt}\selectfont #1}%
    {\fontsize{12pt}{14pt}\selectfont\ #2}
  \end{minipage}
  \par\vspace{0.55em}
}

% Data table header — bold title + italic purpose
\newcommand{\datatable}[2]{%
  \vspace{0.35em}
  {\titlefont\bfseries #1}\ {\titlefont\itshape #2}\par\vspace{0.25em}
}

% Appendix header
\newcommand{\appendixhead}{%
  \par\vspace{0.9em}
  \begin{center}
    {\fontsize{14pt}{16pt}\selectfont Appendix}
  \end{center}
  \vspace{0.4em}
  {\bfseries\fontsize{14pt}{16pt}\selectfont (Scanned data sheets)}
  \par\vspace{0.5em}
}

% ============================================================
% === DOCUMENT BEGINS ===
% ============================================================
\begin{document}

% --- Cover page ---
\includepdf[pages=1,pagecommand={\thispagestyle{empty}}]{\CoverPDF}
\setcounter{page}{1}

% --- Watermark on all content pages ("Physics Labs 2026", plural) ---
\AddToShipoutPictureBG{%
  \begin{tikzpicture}[remember picture,overlay]
    \node[
      rotate=45,
      text=gray!60,
      opacity=0.45,
      scale=5.68,
      font=\ifPDFTeX\normalfont\else\watermarkfont\fi
    ] at ($(current page.center)+(0.25cm,-0.35cm)$) {Physics Labs 2026};
  \end{tikzpicture}%
}

% ============================================================
% 1. ABSTRACT (5 points, ~100 words)
% ============================================================
\labsection{Abstract}{(About 100 words, 5 points)}

% Template:
% This experiment investigates/investigated [topic] using [method].
% [Key quantities] were measured. From [data], [key result with value
% and uncertainty]. The result [agrees with / confirms] [theory/expectation].
% Use \textbf{key terms} in bold.

% Example opening from lab5 (Polarized Light):
% This experiment investigates the \textbf{polarization properties of light}
% using linear polarizers, a quarter-wave plate, and a half-wave plate.
% \textbf{Malus' law} ($I = I_0\cos^2\theta$) was verified by measuring...

\vspace{2cm}  % remove when filled

% ============================================================
% 2. CALCULATIONS, RESULTS, AND COMMENTS (15 points)
%    (or shorter: "Calculations and Results" — both acceptable)
% ============================================================
\labsection{Calculations, Results, and Comments}{(Calculations, data tables and figures; 15 points)}

% --- Structure options ---
% Pattern A: Data tables first → calculations (lab2, lab3)
% Pattern B: Sample calculations first → data tables → error analysis (lab4, lab5)

% --- Sample calculations ---
% Show one complete worked example per formula.
% Box key equations: \boxed{E = (1.54 \pm 0.03) \times 10^{11}\,\mathrm{N\,m^{-2}}}
%
% \textbf{1. Sample Calculations}
% \textit{(a) Descriptive label:}
% \[
% \boxed{\text{key formula}}
% \]
% Substituting measured values: ...

% --- Data tables ---
% \datatable{DATA TABLE X-X}{(Purpose: what this measures)}
% \begin{table}[H]
% \centering
% \small
% \renewcommand{\arraystretch}{1.15}
% \setlength{\tabcolsep}{3pt}
% {\tablefont
% \begin{tabular}{|...|}
% \hline
% ... real data only ...
% \hline
% \end{tabular}
% }
% \caption*{Unit: ... $\Delta_{\mathrm{instr.}}=...$, $\Delta_{\mathrm{read}}=...$}
% \end{table}

% --- Figures (when applicable) ---
% \begin{figure}[H]
% \centering
% \includegraphics[width=0.88\textwidth]{\FigOne}
% \caption{Descriptive caption stating both what is plotted and the key finding.}
% \end{figure}

% --- Error analysis ---
% \textbf{3. Comments and Error Analysis}
% Discussion of results, dominant error sources, comparison to theory.
% List error sources with \textbf{bold labels}:
% \begin{enumerate}[leftmargin=2em,itemsep=0.1em,topsep=0.1em]
% \item \textbf{Angular resolution:} Polarizer scales have $\pm 1^\circ$ precision.
% \item \textbf{Background current:} The 28.5\,\mu\text{A} intercept indicates...
% \end{enumerate}

\vspace{4cm}  % remove when filled

% ============================================================
% 3. CONCLUSIONS (10 points, ~100 words)
% ============================================================
\labsection{Conclusions}{(About 100 words, 10 points)}

% Template:
% [What was verified/measured] was successfully [achieved/verified].
% The final result was \boxed{value ± uncertainty unit}.
% The result is [reasonable / in agreement with theory].
% The main experimental limitation is [dominant error source].

% For multi-topic experiments (like lab5), use \textbf{bold sub-headings}:
% \textbf{Malus' law} was verified: ...
% \textbf{Quarter-wave plate:} The output polarization depends on...
% \textbf{Half-wave plate:} The polarization direction rotates by...

\vspace{3cm}  % remove when filled

% ============================================================
% 4. ANSWERS TO POSTLAB QUESTIONS (10 points)
% ============================================================
\labsection{Answers to Postlab Questions}{(10 points)}

\begingroup
\setstretch{1.02}
\begin{enumerate}[leftmargin=2em,itemsep=0.25em,topsep=0.15em]

\item \textbf{Question title here}

% Best practice (from lab5): bold question summary as first sentence,
% then full explanation with equations, then connection to experimental
% observations.
%
% For questions that ask "how many / what value":
% Start answer with the direct answer: \textbf{Answer: 4 times.}
% Then explain the reasoning.

\item \textbf{Question title here}

% Answer here.

% Add more \item as needed

\end{enumerate}
\endgroup

% ============================================================
% APPENDIX — Scanned Data Sheets
% ============================================================
\clearpage
\appendixhead

Scanned original data sheets used for this report:

\begin{figure}[H]
\centering
\includegraphics[page=1,width=0.90\textwidth]{\DataPDF}
\caption*{Scanned data sheet page 1}
\end{figure}

% For multi-page data sheets:
% \begin{figure}[H]
% \centering
% \includegraphics[page=2,width=0.90\textwidth]{\DataPDF}
% \caption*{Scanned data sheet page 2}
% \end{figure}

% Optional: signature/joke stamp at bottom right
% \vfill
% \hfill\includegraphics[width=2.8cm]{signature.jpg}

\end{document}
